\section{The model: base, experts and likelihood}
\label{sec:model}

\subsection{Why multilayer perceptrons on a characteristic snapshot}

Each input is a point-in-time vector of characteristics for one stock-day, with no time axis. Base and
experts are MLPs on that vector.

\begin{whybox}
\begin{itemize}
\item \textbf{The input format decides the architecture.} Finance mixture-of-experts papers split into those
  using characteristic snapshots (small MLPs: Gu, Kelly \& Xiu, Ye \& Borde) and those using raw price
  sequences (recurrent or transformer experts). The characteristics already summarise each stock's history
  (momentum, volatility, reversal), so a snapshot model is the natural choice.
\item \textbf{Sequence experts would compete with the gate.} An expert with its own memory of recent market
  history could learn the market state itself and absorb exactly the differences between gates the thesis
  tries to measure.
\item \textbf{Benchmark class and feasibility.} Gu, Kelly \& Xiu establish feed-forward networks as the
  benchmark for this problem, and find shallow networks suffice because the signal is weak. Small MLPs make
  thousands of runs feasible on a laptop.
\end{itemize}
\giveup{the possibility that a sequence model finds patterns the characteristics miss.}
\end{whybox}

\subsection{The base network}

\dec{2026-09-21, Q7} The base $f_0$ is an MLP trained on the training block, then frozen (motivation in
Section~1.3). Its own out-of-sample performance is reported beside every result, the headline quantity is
the \emph{improvement over the base}, and its training protocol is fixed in advance.

\begin{whybox}[Why the base protocol is fixed in advance]
An undertrained base leaves general signal for the experts to pick up and makes them look good; an
overtrained one leaves nothing and makes them look useless. Because the base is shared across gates this
does not bias the comparison \emph{between} gates, but it does set the ``does the residual help'' number,
so it must not be tuned after seeing expert results. A null result is possible (the base may capture
most of a weak signal), and how it is reported is decided before results are seen.
\end{whybox}

\subsection{The experts}

Each expert $r_k$ is an MLP on the same 57 inputs, with a \textbf{zero-initialised output layer}, a learned
\textbf{noise scale} $s_k$, and \textbf{pyramid widths} (first hidden width $w$, each further layer half the
previous).

\begin{whybox}
\begin{itemize}
\item \textbf{Same inputs as the base}: an expert must be able to express a regime-specific version of the
  same characteristic-return relation (for example ``reversal is stronger in stress''), so it needs the
  same characteristics.
\item \textbf{Zero initialisation}: the model starts exactly at the base and moves away only where the data
  justify it. It is tested in the closest precedent: random initialisation lost accuracy and made 13 of
  30 seeds collapse.
\item \textbf{A noise scale per expert}: return dispersion is higher in stress, so each regime's expert should be
  allowed its own uncertainty. This is also what the NLL's variance gain measures.
\item \textbf{Pyramid widths}: the convention of Gu, Kelly \& Xiu (e.g.\ 32--16--8). It reduces the width
  choice to one number, which is what makes ``no width sweep'' workable.
\end{itemize}
\end{whybox}

Open details: \open{Q14} whether a shrinkage penalty on the aggregate correction is used (weight
$\alpha$; zero initialisation is tested in the literature, the penalty is not); \open{audit X-1} the
hidden-layer initialisation; \open{Q8} the first width inside the envelope. The load-balancing penalty used
with trained gates is \textbf{inert} under a frozen gate (its gradient is exactly zero), so it is removed or
documented as such.

\subsection{Depth, width and number of experts}

\dec{2026-09-19, Q11} the full gate-by-expert-depth grid; \dec{2026-10-05} the envelope: depths
\textbf{1--3}, first width \textbf{at most 128}, $K$ \textbf{at most 4}; base depth swept standalone.

\begin{whybox}
\begin{itemize}
\item \textbf{Depths 1--3}: Gu, Kelly \& Xiu find performance peaks at three or four layers and their own
  tests cannot separate any pair of depths; deeper adds cost without evidence of benefit. Depths 4--5
  with width 256 would take about a month of laptop time, the only corner compute rules out.
\item \textbf{Base depth standalone}: the base is trained before any gate exists, so its best depth cannot
  depend on a gate, and sweeping it is a clean replication of Gu, Kelly \& Xiu on this panel.
\item \textbf{$K\le4$}: the data cap the number of regimes (next item), and the closest precedent found
  $K=2$ about as good as $K=4$ and $K=6$ worse and less stable.
\end{itemize}
\end{whybox}

\open{Q18} \textbf{$K$}, the number of regimes, is an input to every gate, not an output, and the usable
count is regime \emph{episodes}, not days; over-specifying a mixture slows estimation dramatically
(Heinrich \& Kahn 2018). With the industry-level gate it may become two numbers (Q28). \open{Q17} what
learning theory can say about network size.

\subsection{The likelihood the experts are trained on}

Each expert defines a Gaussian component centred on the base plus its correction:
\begin{equation}
p(y_{i,t}\mid x_{i,t}) \;=\; \sum_{k=1}^{K}\pi_k(i,t)\,\N\!\big(y_{i,t};\,f_0(x_{i,t})+r_k(x_{i,t}),\,s_k^2\big).
\end{equation}
\dec{2026-10-05, Q24} The objective is the \textbf{per-row} likelihood, one log-sum-exp per stock-day,
\begin{equation}
\ell(\psi) \;=\; \sum_{t}\sum_{i\in\mathcal S_t} w_{i,t}\,\log\sum_{k}\pi_k(i,t)\,
\N\!\big(y_{i,t};\,f_0+r_k,\,s_k^2\big),
\end{equation}
with sample weights $w_{i,t}$ (Section~\ref{sec:weights}), read as the \textbf{composite (independence)
likelihood} of the per-date model in which all stocks on a date share one regime.

\begin{whybox}
\begin{itemize}
\item \textbf{The alternative is the per-date likelihood}, where the product over all of about 500 stocks sits inside the
  sum over regimes. It is the ``literal'' model, but it treats 500 correlated returns as independent
  evidence about the regime, so the realised cross-section can overrule the gate almost completely; in
  simulation it becomes distorted whenever regimes are close.
\item \textbf{The per-row form is consistent}: each term is the exact one-stock marginal of the per-date model,
  so its score has mean zero and the estimates converge to the right values even though cross-stock
  dependence is ignored (composite likelihood theory; Varin, Reid \& Firth 2011).
\item \textbf{It is robust} when the independence assumption fails, which for stock returns it partly does.
\item \textbf{It keeps the gate in charge}, which is the point of freezing it.
\item \textbf{Inference does not rely on it}: significance comes from walk-forward rank ICs with
  overlap-corrected standard errors, so its understated standard errors do no harm.
\end{itemize}
With the industry-level gate the latent is per industry and date, and the per-row form is closer to exact.
\giveup{statistical efficiency (the full joint likelihood would use the dependence between stocks if it
were correctly specified).}
\end{whybox}

\open{Q20} Whether the mixture likelihood stays the training objective at all, against squared or Huber
error on the point forecast or a rank/IC objective. Under the likelihood each expert's gradient is
weighted by its posterior responsibility, which depends on the realised target, so the outcome partly
re-assigns experts after the frozen gate has assigned them; under squared error the gate's weight alone
decides.

\subsection{The point forecast and the NLL}

The forecast used for ranking and portfolios is the mixture mean, equation~\eqref{eq:model}. The NLL is
reported three ways (\texttt{NLL\_single}, \texttt{NLL\_base}, \texttt{NLL\_full}), with the
volatility-only gain next to the correction gain, so that a gain from the experts' different noise
scales is never mistaken for better return forecasts.

\subsection{Planned construction arms}

A construction pilot (one gate held fixed) with five arms: (A) standard MoE, joint training; (B) residual
MoE with the base pre-trained and then unfrozen; (C) residual MoE, frozen base (the design); (D) C with
randomly initialised experts; (E) C with $\alpha=0$. \open{Q6 design half, Q14}

\begin{whybox}[Why run them]
They test the construction instead of assuming it. A versus C checks that the residual advantage found
for volatility survives on returns; B versus C isolates the effect of freezing the base, the one cell
the literature has not tested; D and E separate zero initialisation from the shrinkage penalty.
\end{whybox}
